\section{Introduction}
  \label{sec:introduction}

  Gravity is traditionally described as a geometric phenomenon.
  In general relativity, the gravitational interaction is encoded in the curvature of spacetime, and the Newtonian
  potential emerges as a weak-field limit of the Einstein equations.
  This geometric formulation has been extraordinarily successful, yet it leaves open the conceptual question of whether
  gravitational interaction must be fundamental, or whether it could arise as an effective response of a more primitive
  structure.

  In parallel, operator-theoretic and spectral viewpoints have increasingly clarified the role of elliptic operators in
  encoding geometric and physical information.
  In particular, functional determinants and resolvents of Laplace-type operators play a central role in quantum field
  theory, statistical mechanics, and spectral geometry.
  Variations of log-determinants are known to generate non-local responses through the inverse operator, suggesting a
  possible bridge between entropy-like functionals and long-range interactions.

  The purpose of this work is to explore a minimal mechanism by which a Newtonian interaction may emerge from the
  variation of a projective entropy functional.
  We consider a positive elliptic operator $A$
  acting on a three-dimensional domain and define an entropy functional proportional to the logarithm of its (
  regularized) determinant.
  We show that, in the weak-perturbation regime, the variation of this functional reduces to a resolvent trace,
  \begin{equation}
    \delta S_\Pi \sim \tfrac12 \mathrm{Tr}(A^{-1} \delta A),
  \end{equation}
  so that the Green kernel of $A$ governs the spatial structure of the response.

  Since the Green function of a Laplace-type operator in three dimensions decays as $1/r$
  , the resulting entropy variation inherits a Newtonian kernel.
  This mechanism does not assume a fundamental force law.
  Instead, long-range interaction arises as a non-local consequence of the resolvent structure of the operator.

  We further demonstrate numerically that the emergence of a macroscopic $1/r$
  amplitude shift depends on the spatial support of the perturbation.
  Compact perturbations confined to a sufficiently small region modify the entropy locally without generating a
  measurable long-range effect.
  By contrast, perturbations extending over a finite spatial region induce a coherent modification of the Green
  amplitude consistent with a Newtonian response.

  The results presented here provide operator-theoretic evidence that a gravitational interaction can arise as a
  resolvent-mediated effect of projective entropy variation.
  Functional determinants and spectral invariants play a central role
  in quantum geometry and effective field theory~\cite{dowker1976,polyakov1981}.
  The analysis is intentionally minimal and does not rely on additional geometric assumptions.
  Its scope is restricted to the Newtonian regime, leaving relativistic extensions for future investigation.
